% !TEX root = ../../../main.tex
% Appendix entry: catalogue.cnn_score -- "Model stage (CNN)" (Appendix A).
\entryfield{Algorithm}
\code{_run} and \code{score_stack} in \code{Adapters/catalogue_cnn_score.py}, over the image
stack from \blockref{catalogue.window_images} and the \code{windows} side input:
\begin{enumerate}
\item Refuse an input that is not an image \code{Encoding}, a missing side input, a stack not
  shaped $(N, H, W, 3)$, or a stack whose $N$ differs from the window count.
\item \emph{Load} (\code{load_model} in \code{Working/Catalogue/cnn/apply_cnn.py}): read the
  state dict at \code{model_path}, take the class count from the shape of
  \code{backbone.classifier.1.weight}, build \code{EEG_CNN(num_classes)} and load the weights;
  evaluation mode, on CUDA if available, else the CPU.
\item \emph{Transform} each image with \code{get_transforms(img_size, is_train=False,
  is_rgb=True)} from \code{Working/Catalogue/cnn/cnn_rangapur.py}: resize to
  \code{img_size} $\times$ \code{img_size}, convert to a tensor in $[0,1]$, normalise each
  channel with mean 0.5 and standard deviation 0.5.
\item \emph{Infer} in batches of \code{batch_size} without gradients;
  $p_i = \mathrm{softmax}(\text{logits}_i)_0$, the probability of class 0.
\item \emph{Lay out}: a score array of the span's length filled with NaN; for each window
  (start minus $\mathrm{round}(t_0 f_s)$) inside the span, its $L$ samples are set to $p_i$.
  Where windows overlap the later window wins.
\end{enumerate}
\code{EEG_CNN} is an EfficientNet-B0 \citep{tan2019efficientnet} from torchvision with its
head replaced by dropout ($p = 0.3$) and a linear layer to \code{num_classes} outputs. Class 0
is ``interesting'' because torchvision's \code{ImageFolder} numbers the class folders
alphabetically (\code{interesting} before \code{notinteresting}).

\entryfield{Assumptions}
\begin{itemize}
\item \emph{The checkpoint.} Two classes, class 0 interesting, trained on the image type the
  stack holds. Since 2026-10-17 (W4-A) the class count is checked: a checkpoint with other than two
  classes is refused, naming the file and the count; the image type and the class order are still
  not checked. \code{docs/prompts/fixup/reports/AH-blind-test-labelling.md}
  records, for the checkpoints in \code{MODELS/}: \code{fusion_cnn.pth} has \emph{one}
  output class, so its softmax is 1.0 for every window; \code{fusion_cnn_2.pth} and
  \code{fusion_cnn_3.pth} have two, with class 0 interesting (AUC 0.979 and 0.970 on 120
  hand-labelled 10-minute M2\_aug windows, very likely part of their training data);
  \code{fusion_prediction_cnn.pth} has 14 classes and answers a different question.
\item \emph{Training.} \code{cnn_rangapur.py} is the training script: ImageNet-pretrained
  EfficientNet-B0, a random 80/20 train/validation split of the image folders, Adam at
  $10^{-3}$ with cosine decay, up to 25 epochs with early stopping on validation accuracy,
  random horizontal and vertical flips and colour jitter, on 10-minute windows at 1\,Hz.
  \unverified{That the shipped checkpoints were made with these settings} is not recorded.
\item \emph{Window geometry.} Starts are channel-absolute; $t_0$ is the span's first time in
  seconds from the channel's start.
\item \emph{Units and rate.} The images are unit-free (see \blockentry{catalogue.window_images});
  the output \code{Scores} carry the span's $f_s$.
\end{itemize}

\entryfield{Outputs}
\code{Scores} over the span: $p \in [0,1]$ per sample, NaN where no window lies. \code{meta}
holds \code{model_path}, \code{n_windows}, \code{window_scores} and \code{starts} (the first
4{,}096 of each), \code{length} and \code{p_range}.

\entryfield{Cost}
One forward pass of EfficientNet-B0 per window: linear in $N$, much faster on a GPU. The
model is loaded on every run. \code{estimate} reads the calibration file (calibrated by
scoring random images with the default checkpoint) and returns \code{None} until calibrated.
On the cluster: a recipe exported for SLURM that contains this block is sent to the GPU
partition with an A100 and CUDA (\code{recipe_uses_gpu} in \code{Working/hpc/job_export.py},
whose \code{_GPU_ALGORITHMS} names \code{cnn_score}); every other block in this group runs on
the CPU partition. The block itself runs locally too, on the CPU when there is no GPU.

\entryfield{Failure modes}
\begin{itemize}
\item \emph{The default checkpoint is refused.} \code{model_path} defaults to
  \code{MODELS/fusion_cnn.pth}, the one-class network, which would score every window 1.0. Since
  2026-10-17 the block raises a ValueError naming the file and its one class (the default was
  left unchanged so no recipe hash moved); pick \code{fusion_cnn_3.pth}. The same check refuses
  \code{fusion_prediction_cnn.pth} (14 classes).
\item \emph{Image type mismatch.} A GASF network fed fusion pictures, or the reverse, scores
  without complaint.
\item \emph{Pretrained weights} are no longer fetched when a checkpoint is loaded:
  \code{apply_cnn.load_model} builds \code{EEG_CNN(pretrained=False)} since 2026-10-17, so an
  offline machine scores; training still starts from ImageNet weights.
\item \emph{Training augmentations.} Single-axis flips of a Gramian or recurrence image are
  not images of any signal (\code{docs/prompts/fixup/reports/AI-cnn-arm-slurm.md}); the
  networks were trained with them.
\item \emph{Contamination.} Scores on the hand-labelled windows are an upper bound, not a
  measure of generalisation.
\item \emph{Partial windows at the span's end} are written as far as the span goes; windows
  starting before the span are skipped.
\end{itemize}

\entryfield{Validation}
\code{tests/test_adapter_cnn_score.py} checks the plumbing: the declared types, the
refusals, and that each window's samples carry one value in $[0,1]$, with a random-weight
two-class checkpoint written to a temporary directory; it does not check the answer.
\code{tests/test_cnn_window_scores.py} checks that this block and
\blockref{catalogue.cnn_window_scores} give the same score on the same checkpoint, and
\code{tests/test_adapter_cnn_window_scores.py} that a one-class checkpoint is refused. The
AUCs above come from a separate check recorded in the AH report, not from a test.

\entryfield{References}
\citet{tan2019efficientnet}; the images, \citet{wang2015imaging} and
\citet{eckmann1987recurrence}.

\entryfield{Tests}
\code{tests/test_adapter_cnn_score.py}. \code{tests/test_block_standard.py} checks that it
declares a cost estimate. \code{tests/test_job_export.py} checks that a window matrix with a
CNN stage is sent to the GPU partition; no test checks the same for a recipe containing this
block.
